\documentclass[10pt,a4paper]{amsart}
\usepackage{amsmath,amssymb,amsthm,graphicx}
\usepackage{verbatim}
\newenvironment{reference}{\comment}{\endcomment}
%\newenvironment{reference}{}{}
\newenvironment{slogan}{\comment}{\endcomment}
\newenvironment{history}{\comment}{\endcomment}

% For commutative diagrams we use Xy-pic
\usepackage[all]{xy}

% We use 2cell for 2-commutative diagrams.
\xyoption{2cell}
% Theorem environments.
%
\theoremstyle{plain}
\newtheorem{theorem}[subsection]{Theorem}
\newtheorem{proposition}[subsection]{Proposition}
\newtheorem{lemma}[subsection]{Lemma}

\theoremstyle{definition}
\newtheorem{definition}[subsection]{Definition}
\newtheorem{example}[subsection]{Example}
\newtheorem{exercise}[subsection]{Exercise}
\newtheorem{situation}[subsection]{Situation}

\theoremstyle{remark}
\newtheorem{remark}[subsection]{Remark}
\newtheorem{remarks}[subsection]{Remarks}

\numberwithin{equation}{subsection}

% Macros
%
\def\lim{\mathop{\mathrm{lim}}\nolimits}
\def\colim{\mathop{\mathrm{colim}}\nolimits}
\def\Spec{\mathop{\mathrm{Spec}}}
\def\Hom{\mathop{\mathrm{Hom}}\nolimits}
\def\Ext{\mathop{\mathrm{Ext}}\nolimits}
\def\SheafHom{\mathop{\mathcal{H}\!\mathit{om}}\nolimits}
\def\SheafExt{\mathop{\mathcal{E}\!\mathit{xt}}\nolimits}
\def\Sch{\mathit{Sch}}
\def\Mor{\mathop{\mathrm{Mor}}\nolimits}
\def\Ob{\mathop{\mathrm{Ob}}\nolimits}
\def\Sh{\mathop{\mathit{Sh}}\nolimits}
\def\NL{\mathop{N\!L}\nolimits}
\def\CH{\mathop{\mathrm{CH}}\nolimits}
\def\proetale{{pro\text{-}\acute{e}tale}}
\def\etale{{\acute{e}tale}}
\def\QCoh{\mathit{QCoh}}
\def\Ker{\mathop{\mathrm{Ker}}}
\def\Im{\mathop{\mathrm{Im}}}
\def\Coker{\mathop{\mathrm{Coker}}}
\def\Coim{\mathop{\mathrm{Coim}}}

% Boxtimes
%
\DeclareMathSymbol{\boxtimes}{\mathbin}{AMSa}{"02}

%
% Macros for moduli stacks/spaces
%
\def\QCohstack{\mathcal{QC}\!\mathit{oh}}
\def\Cohstack{\mathcal{C}\!\mathit{oh}}
\def\Spacesstack{\mathcal{S}\!\mathit{paces}}
\def\Quotfunctor{\mathrm{Quot}}
\def\Hilbfunctor{\mathrm{Hilb}}
\def\Curvesstack{\mathcal{C}\!\mathit{urves}}
\def\Polarizedstack{\mathcal{P}\!\mathit{olarized}}
\def\Complexesstack{\mathcal{C}\!\mathit{omplexes}}
% \Pic is the operator that assigns to X its picard group, usage \Pic(X)
% \Picardstack_{X/B} denotes the Picard stack of X over B
% \Picardfunctor_{X/B} denotes the Picard functor of X over B
\def\Pic{\mathop{\mathrm{Pic}}\nolimits}
\def\Picardstack{\mathcal{P}\!\mathit{ic}}
\def\Picardfunctor{\mathrm{Pic}}
\def\Deformationcategory{\mathcal{D}\!\mathit{ef}}
\usepackage{cleveref}
\crefname{originalreference}{original reference}{original references}
\Crefname{originalreference}{Original reference}{Original references}
\makeatletter
\newlabel{dga-lemma-tensor-hom-adjunction}{{09LN}{source}{}{sourcepointer-1}{}}
\newlabel{dga-lemma-tensor-hom-adjunction@cref}{{[originalreference][0][]09LN}{[1][0][]source}{}{}{}}
\makeatother
\crefname{originalreference}{original reference}{original references}
\Crefname{originalreference}{Original reference}{Original references}
\makeatletter
\newlabel{categories-lemma-adjoint-exact}{{0038}{source}{}{sourcepointer-2}{}}
\newlabel{categories-lemma-adjoint-exact@cref}{{[originalreference][0][]0038}{[1][0][]source}{}{}{}}
\makeatother
\crefname{originalreference}{original reference}{original references}
\Crefname{originalreference}{Original reference}{Original references}
\makeatletter
\newlabel{homology-section-functors}{{010M}{source}{}{sourcepointer-3}{}}
\newlabel{homology-section-functors@cref}{{[originalreference][0][]010M}{[1][0][]source}{}{}{}}
\makeatother
\title{Four characterizations of functors Mod\_A → Mod\_B: tensoring, right exact + direct sums, all colimits, right adjoint}
\author{The Stacks Project Authors}
\date{stacks-project:functors.tex:local-snapshot-20261012; GNU FDL 1.2 or later}
\begin{document}
\maketitle
\noindent\textit{Editorial note.} All delivered mathematical text is the author's own, copied verbatim from the pinned Stacks Project chapter file functors.tex; no statement or proof step was rewritten, reordered or removed.\par\medskip
\noindent\textit{Editorial note.} Span recomputation by label pairing (LANE\_BRIEF 10.2): in this chapter the author's proof environment follows the statement environment as a standalone block; the spec covers the complete environment pair (statement span, proof span) as recorded in the pool candidate.\par\medskip
\noindent\textit{Editorial note.} Occupation scan: combined claim\_key + statement\_anchor + generic-line scan over all registration files (active and withdrawn) found 0 hits for functors.tex (fu-occupied.json, 2026-10-14); consistent with the duty manager's six-check dispatch.\par\medskip
\ifdefined\hypertarget\hypertarget{sourcepointer-1}{}\fi
\noindent\textit{Source reference.} Original reference 09LN: sources/stacks-chow/tags.txt, source lines 12151--12151. Cross-chapter standard prerequisite cited by the author inside the counted span; it is used as a pointer and its own proof is not part of this item.\par\medskip
\ifdefined\hypertarget\hypertarget{sourcepointer-2}{}\fi
\noindent\textit{Source reference.} Original reference 0038: sources/stacks-chow/tags.txt, source lines 125--125. Cross-chapter standard prerequisite cited by the author inside the counted span; it is used as a pointer and its own proof is not part of this item.\par\medskip
\ifdefined\hypertarget\hypertarget{sourcepointer-3}{}\fi
\noindent\textit{Source reference.} Original reference 010M: sources/stacks-chow/tags.txt, source lines 1332--1332. Cross-chapter standard prerequisite cited by the author inside the counted span; it is used as a pointer and its own proof is not part of this item.\par\medskip
\begin{lemma}
\label{lemma-functor}
Let $A$ and $B$ be rings. Let $F : \text{Mod}_A \to \text{Mod}_B$
be a functor. The following are equivalent
\begin{enumerate}
\item $F$ is isomorphic to the functor $M \mapsto M \otimes_A K$
for some $A \otimes_\mathbf{Z} B$-module $K$,
\item $F$ is right exact and commutes with all direct sums,
\item $F$ commutes with all colimits,
\item $F$ has a right adjoint $G$.
\end{enumerate}
\end{lemma}

\begin{proof}
If (1), then (4) as a right adjoint for $M \mapsto M \otimes_A K$
is $N \mapsto \Hom_B(K, N)$, see
Differential Graded Algebra, Lemma \ref{dga-lemma-tensor-hom-adjunction}.
If (4), then (3) by Categories, Lemma \ref{categories-lemma-adjoint-exact}.
The implication (3) $\Rightarrow$ (2) is immediate from the definitions.

\medskip\noindent
Assume (2). We will prove (1). By the discussion in
Homology, Section \ref{homology-section-functors}
the functor $F$ is additive. Hence $F$ induces
a ring map $A \to \text{End}_B(F(M))$, $a \mapsto F(a \cdot \text{id}_M)$
for every $A$-module $M$. We conclude that $F(M)$ is an
$A \otimes_\mathbf{Z} B$-module functorially in $M$.
Set $K = F(A)$. Define
$$
M \otimes_A K = M \otimes_A F(A) \longrightarrow F(M),
\quad m \otimes k \longmapsto F(\varphi_m)(k)
$$
Here $\varphi_m : A \to M$ sends $a \to am$. The rule
$(m, k) \mapsto F(\varphi_m)(k)$ is $A$-bilinear (and $B$-linear
on the right) as required to obtain the displayed
$A \otimes_\mathbf{Z} B$-linear map.
This construction is functorial in $M$, hence defines a transformation
of functors $- \otimes_A K \to F(-)$ which is an isomorphism when
evaluated on $A$. For every $A$-module $M$ we can choose an exact sequence
$$
\bigoplus\nolimits_{j \in J} A \to
\bigoplus\nolimits_{i \in I} A \to
M \to 0
$$
Using the maps constructed above we find a commutative diagram
$$
\xymatrix{
(\bigoplus\nolimits_{j \in J} A) \otimes_A K \ar[r] \ar[d] &
(\bigoplus\nolimits_{i \in I} A) \otimes_A K \ar[r] \ar[d] &
M \otimes_A K \ar[r] \ar[d] &
0 \\
F(\bigoplus\nolimits_{j \in J} A) \ar[r] &
F(\bigoplus\nolimits_{i \in I} A) \ar[r] &
F(M) \ar[r] &  0
}
$$
The lower row is exact as $F$ is right exact.
The upper row is exact as tensor product with $K$ is right exact.
Since $F$ commutes with direct sums the left two vertical arrows
are bijections. Hence we conclude.
\end{proof}

\end{document}
